\documentclass[11pt,a4paper]{article}
\usepackage[T1]{fontenc}
\usepackage[utf8]{inputenc}
\usepackage[french]{babel}
\usepackage{lmodern}
\usepackage{microtype}
\usepackage[margin=2.2cm]{geometry}
\usepackage{amsmath,amssymb,amsthm,mathtools}
\usepackage{icomma}
\usepackage{booktabs,tabularx,array}
\usepackage{xcolor}
\usepackage{tikz}
\usetikzlibrary{patterns,positioning}
\usepackage{pgfplots}
\pgfplotsset{compat=1.18}
\usepackage[most]{tcolorbox}
\usepackage{enumitem}
\usepackage[hidelinks]{hyperref}
\usepackage{caption}
\usepackage{float}

\definecolor{nyt}{HTML}{1F3B63}
\definecolor{wp}{HTML}{17695B}
\definecolor{gov}{HTML}{B42318}
\definecolor{nat}{HTML}{9A6212}
\definecolor{regnone}{HTML}{C6C2B8}
\definecolor{regnosuit}{HTML}{CBE2C3}
\definecolor{regnyt}{HTML}{F0C5BD}
\definecolor{regboth}{HTML}{F1D69B}

\newtcolorbox{cours}{colback=nat!6,colframe=nat!70,boxrule=0.5pt,arc=2pt,left=6pt,right=6pt,top=4pt,bottom=4pt,fonttitle=\bfseries\small,title=Lien avec le cours}
\newtcolorbox{resultat}{colback=wp!6,colframe=wp!70,boxrule=0.5pt,arc=2pt,left=6pt,right=6pt,top=4pt,bottom=4pt}
\newtcolorbox{intuition}{colback=nyt!4,colframe=nyt!50,boxrule=0.4pt,arc=2pt,left=6pt,right=6pt,top=3pt,bottom=3pt,fonttitle=\bfseries\small,title=Intuition}

\theoremstyle{definition}
\newtheorem{lemme}{Lemme}
\newtheorem{proposition}{Proposition}
\newcommand{\ind}[1]{\mathbf{1}_{\{\text{#1}\}}}
\newcommand{\E}{\mathbb{E}}
\setlength{\parindent}{0pt}
\setlength{\parskip}{5pt}
\captionsetup{font=small,labelfont=bf}

\title{\textbf{Pentagon Papers}\\[3pt]\large Modèle, équations et calcul des probabilités de favorabilité des juges}
\author{Travail de théorie des jeux, groupe 5--6 · ECGEB366}
\date{}

\begin{document}
\maketitle
\vspace{-1.2cm}

\begin{resultat}
\textbf{En bref.} Trois joueurs (le Times, le Post, l'administration Nixon) jouent un jeu séquentiel à information parfaite. Les juges ne sont pas des joueurs : leurs verdicts sont incertains et entrent dans les gains par deux probabilités,
\begin{align*}
q&=\Pr(\text{un juge local refuse la censure}), & p&=\Pr(\text{la Cour suprême donne raison aux journaux}).
\end{align*}
Avec $q=0,3$ et $p=0,65$, l'équilibre de Nash parfait en sous-jeux (ENPS) reproduit 1971 : le Times publie, Nixon le poursuit, le Post prend le relais, Nixon le poursuit aussi. Les seuils clés sont :
\begin{itemize}[leftmargin=*,itemsep=0pt,topsep=2pt]
\item le Post publie en relais si $p\ge\frac{53}{88}\approx60,2\,\%$ ;
\item Nixon poursuit le Post si $p<\frac67\approx85,7\,\%$ ;
\item Nixon poursuit le Times si $q<\frac{19}{29}\approx65,5\,\%$ (avec $p=0,65$).
\end{itemize}
Avant le premier article, la probabilité que tout sorte vaut $q+(1-q)p=75,5\,\%$, celle d'une censure confirmée $(1-q)(1-p)=24,5\,\%$.
\end{resultat}

\section{Le jeu}

\subsection{Joueurs, actions et ordre}
\begin{itemize}[leftmargin=*,itemsep=1pt]
\item \textbf{Le New York Times} ($T$) décide de publier ou non (nœud $x_0$).
\item \textbf{L'administration Nixon} ($N$) poursuit ou laisse faire, à trois nœuds : face au Times ($x_1$), face au Post en relais ($x_3$), face au Post seul ($x_5$).
\item \textbf{Le Washington Post} ($W$) publie ou non, à deux nœuds : en relais du Times ($x_2$) ou seul si le Times renonce ($x_4$).
\end{itemize}
Les \textbf{juges} n'ont ni gains ni choix stratégique : ce ne sont pas des joueurs, ils n'ont pas de nœud. Ellsberg remet les documents au Times, puis au Post (coup fixé).

\begin{cours}
Un jeu se décrit par ses joueurs, leurs stratégies et leurs gains (S5 ; S6, diapositives 7--9). Information parfaite : chaque ensemble d'information contient un seul nœud (S6, diapositives 7 et 36--39). Une stratégie est un plan complet : le Times en a $2$, le Post $2^2=4$, Nixon $2^3=8$, d'où $2\times4\times8=64$ profils (S6, diapositives 14--18).
\end{cours}

\subsection{L'arbre}
\begin{figure}[H]
\centering
\begin{tikzpicture}[x=2.35cm,y=1.55cm,
  pl/.style={circle,inner sep=2.6pt,fill=#1},
  lab/.style={font=\scriptsize,inner sep=1.2pt,fill=white},
  eq/.style={line width=1.6pt,draw=gov},
  br/.style={line width=0.6pt},
  leaf/.style={font=\scriptsize,align=center,inner sep=1pt}]
% positions : feuilles en colonnes 0..6, nœuds au milieu
\coordinate (x0) at (3.6875,0);
\coordinate (x1) at (2.125,-1);
\coordinate (x2) at (1.25,-2);
\coordinate (x3) at (0.5,-3);
\coordinate (x4) at (5.25,-1);
\coordinate (x5) at (4.5,-2);
\coordinate (O1) at (0,-4); \coordinate (O2) at (1,-4); \coordinate (O3) at (2,-3);
\coordinate (O4) at (3,-2); \coordinate (O5) at (4,-3); \coordinate (O6) at (5,-3); \coordinate (O7) at (6,-2);
% branches (traits épais : ENPS)
\draw[eq] (x0)--(x1) node[lab,midway,above left]{Publier};
\draw[br,dashed] (x0)--(x4) node[lab,midway,above right]{Ne pas publier};
\draw[eq] (x1)--(x2) node[lab,midway,left]{Poursuivre};
\draw[br,dashed] (x1)--(O4) node[lab,midway,right]{Laisser faire};
\draw[eq] (x2)--(x3) node[lab,midway,left]{Publier};
\draw[br,dashed] (x2)--(O3) node[lab,midway,right]{Ne pas publier};
\draw[eq] (x3)--(O1) node[lab,midway,left]{Poursuivre};
\draw[br,dashed] (x3)--(O2) node[lab,midway,right]{Laisser faire};
\draw[br,dashed] (x4)--(x5) node[lab,midway,left]{Publier};
\draw[eq] (x4)--(O7) node[lab,midway,right]{Ne pas publier};
\draw[eq] (x5)--(O5) node[lab,midway,left]{Poursuivre};
\draw[br,dashed] (x5)--(O6) node[lab,midway,right]{Laisser faire};
% nœuds des joueurs
\node[pl=nyt,label={[font=\small\bfseries,text=nyt]above:Times $x_0$}] at (x0) {};
\node[pl=gov,label={[font=\small\bfseries,text=gov]above left:Nixon $x_1$}] at (x1) {};
\node[pl=wp,label={[font=\small\bfseries,text=wp]above left:Post $x_2$}] at (x2) {};
\node[pl=gov,label={[font=\small\bfseries,text=gov]above left:Nixon $x_3$}] at (x3) {};
\node[pl=wp,label={[font=\small\bfseries,text=wp]above right:Post $x_4$}] at (x4) {};
\node[pl=gov,label={[font=\small\bfseries,text=gov]above left:Nixon $x_5$}] at (x5) {};
% feuilles : gains (Times ; Post ; Nixon), espérés si procès
\foreach \n/\t/\g in {O1/{Procès du Post\\+ tribunaux}/{60,6\\8,82\\$-$189,7},O2/{Post publie,\\pas de procès}/{90\\84\\$-$210},O3/{Times censuré,\\Post muet}/{6\\0\\$-$112},O4/{Tout sort,\\sans procès}/{100\\0\\$-$200},O5/{Procès du\\Post seul}/{$-$25,5\\$-$79,8\\$-$122},O6/{Le Post\\publie seul}/{$-$50\\120\\$-$200},O7/{Rien\\ne sort}/{0\\0\\0}}{
  \fill (\n) circle (1.6pt);
  \node[leaf,below=2pt] at (\n) {\textbf{\n}\\[-1pt]\t\\[1pt]\g};
}
\end{tikzpicture}
\caption{Forme étendue avec les trois joueurs seulement. Gains dans l'ordre (Times ; Post ; Nixon). Les feuilles $O_1,O_2,O_3,O_5$ portent des gains \emph{espérés} qui intègrent les verdicts des juges. Traits épais : actions de l'ENPS ; pointillés : actions écartées.}
\end{figure}

\subsection{Paramètres}
\begin{center}
\small
\begin{tabular}{@{}lllp{8.2cm}@{}}
\toprule
Symbole & Valeur & Joueur & Signification \\
\midrule
$q$ & $0,3$ & juges & probabilité qu'un juge local refuse la censure \\
$p$ & $0,65$ & juges & probabilité que la Cour suprême donne raison aux journaux \\
$\theta$ & $0,3$ & tous & part des révélations publiée avant une censure (trois articles) \\
$B_T,\ B_W$ & $100,\ 120$ & Times, Post & valeur d'une publication complète (scoop, prestige, mission) \\
$c_T,\ c_W$ & $10,\ 15$ & Times, Post & frais totaux d'un procès \\
$S_T,\ S_W$ & $50,\ 180$ & Times, Post & sanction si la censure est maintenue (bourse et licences pour le Post) \\
$R_T$ & $50$ & Times & perte si le Post publie le scoop que le Times a refusé \\
$D$ & $200$ & Nixon & dommage si tout sort \\
$k_T,\ k_W$ & $10,\ 20$ & Nixon & coût d'un procès contre le Times, contre le Post \\
\bottomrule
\end{tabular}
\end{center}

\section{Les gains}

\subsection{Les fonctions de gain}
Soit $\sigma$ la part des documents qui finit par être publique : $\sigma=1$ si tout sort, $\sigma=\theta$ si une censure est maintenue, $\sigma=0$ si rien ne sort. Avec l'indicatrice $\mathbf{1}_A$ ($1$ si $A$ se produit, $0$ sinon) :
\begin{align}
u_T &= \ind{Times publie}\,B_T\,\sigma-\ind{Times poursuivi}\,c_T \notag\\
&\quad-\ind{censure du Times maintenue}\,S_T-\ind{Times renonce, Post publie}\,R_T\,\sigma, \label{eq:uT}\\
u_W &= \ind{Post publie}\,B_W\,\sigma-\ind{Post poursuivi}\,c_W-\ind{censure du Post maintenue}\,S_W, \label{eq:uW}\\
u_N &= -D\,\sigma-\ind{Times poursuivi}\,k_T-\ind{Post poursuivi}\,k_W. \label{eq:uN}
\end{align}

\begin{intuition}
Un journal gagne en proportion de ce qu'il réussit à rendre public et perd ce que le procès lui coûte ; une censure maintenue ajoute une sanction. Nixon perd en proportion de ce qui sort et paie chaque procès. Le zéro est « rien ne se passe » : ajouter une constante aux gains d'un joueur ne change aucun de ses choix.
\end{intuition}

\subsection{Les situations finales}
Une fois les juges passés, il ne reste que dix situations possibles. Chacune s'obtient en appliquant \eqref{eq:uT}--\eqref{eq:uN}.
\begin{center}
\footnotesize
\begin{tabularx}{\textwidth}{@{}>{\raggedright\arraybackslash}Xcccc@{}}
\toprule
Situation finale & $\sigma$ & $u_T$ & $u_W$ & $u_N$ \\
\midrule
Le juge local refuse de censurer le Times & $1$ & $B_T-c_T=90$ & $0$ & $-D-k_T=-210$ \\
Cour suprême : victoire des journaux & $1$ & $B_T-c_T=90$ & $B_W-c_W=105$ & $-D-k_T-k_W=-230$ \\
Cour suprême : censure confirmée & $\theta$ & $\theta B_T-c_T-S_T=-30$ & $\theta B_W-c_W-S_W=-159$ & $-\theta D-k_T-k_W=-90$ \\
Times censuré, le Post publie librement & $1$ & $B_T-c_T=90$ & $B_W=120$ & $-D-k_T=-210$ \\
Times censuré, le Post se tait & $\theta$ & $\theta B_T-c_T-S_T=-30$ & $0$ & $-\theta D-k_T=-70$ \\
Le Times publie, aucun procès & $1$ & $B_T=100$ & $0$ & $-D=-200$ \\
Le juge refuse de censurer le Post seul & $1$ & $-R_T=-50$ & $B_W-c_W=105$ & $-D-k_W=-220$ \\
Le Post seul est censuré & $\theta$ & $-\theta R_T=-15$ & $\theta B_W-c_W-S_W=-159$ & $-\theta D-k_W=-80$ \\
Le Post publie seul, aucun procès & $1$ & $-R_T=-50$ & $B_W=120$ & $-D=-200$ \\
Personne ne publie & $0$ & $0$ & $0$ & $0$ \\
\bottomrule
\end{tabularx}
\end{center}

\begin{cours}
Un vecteur de gains à chaque feuille, dans un ordre fixe (S6, diapositives 8--9 et 14). Une feuille résume une histoire complète : « tout sort » ne vaut pas la même chose avec ou sans procès (S6, diapositives 14--18). Un sous-jeu reprend les gains du jeu entier : les mêmes règles s'appliquent partout (S6, diapositives 10--13).
\end{cours}

\section{Les juges : des probabilités, pas des joueurs}

\subsection{Les verdicts d'un procès}
Un procès commence devant un juge local. S'il refuse la censure (probabilité $q$), tout sort. S'il censure ($1-q$) et que l'autre journal publie aussi et est poursuivi, les deux affaires sont jointes devant la Cour suprême, qui donne raison aux journaux avec la probabilité $p$. On suppose les deux verdicts indépendants et les probabilités connues de tous. Dans le procès du Post après celui du Times, les trois verdicts possibles ont pour probabilités
\[
\underbrace{q}_{\text{refus local}}\;+\;\underbrace{(1-q)\,p}_{\text{censure, puis Cour favorable}}\;+\;\underbrace{(1-q)(1-p)}_{\text{censure confirmée}}\;=\;q+(1-q)\big[p+(1-p)\big]\;=\;1 .
\]
La probabilité que la justice donne \emph{finalement} raison aux journaux vaut donc
\begin{equation}
\pi \;=\; q+(1-q)\,p \;=\; 1-(1-q)(1-p) \;=\; 0,3+0,7\times0,65 \;=\; 0,755. \label{eq:pi}
\end{equation}
Sans second journal, la Cour suprême n'est pas saisie en urgence : un journal seul n'a que le juge local, avec $\pi=q$.

\subsection{Les gains espérés des feuilles}
Chaque feuille avec procès est une loterie sur les situations finales. Son gain est l'espérance :
\begin{equation}
\E\big[u_i(O_k)\big]=\sum_{v}\Pr(v)\,u_i\big(\text{situation après le verdict } v\big),\qquad i\in\{T,W,N\}. \label{eq:esp}
\end{equation}

\paragraph{$O_1$ — procès du Post + tribunaux.}
\begin{align*}
\E[u_T(O_1)] &= q(B_T-c_T)+(1-q)\big[p(B_T-c_T)+(1-p)(\theta B_T-c_T-S_T)\big]\\
&= 0,3\times90+0,7\,\big[0,65\times90+0,35\times(-30)\big] = 27+0,7\times48 = 60,6,\\[3pt]
\E[u_W(O_1)] &= q\times0+(1-q)\big[p(B_W-c_W)+(1-p)(\theta B_W-c_W-S_W)\big]\\
&= 0,7\,\big[0,65\times105+0,35\times(-159)\big] = 0,7\times12,6 = 8,82,\\[3pt]
\E[u_N(O_1)] &= q(-D-k_T)+(1-q)\big[p(-D-k_T-k_W)+(1-p)(-\theta D-k_T-k_W)\big]\\
&= 0,3\times(-210)+0,7\,\big[0,65\times(-230)+0,35\times(-90)\big] = -63+0,7\times(-181) = -189,7.
\end{align*}
\paragraph{$O_2$ — le Post publie, pas de procès contre lui.}
$\E[u_T]=B_T-c_T=90$, \quad $\E[u_W]=(1-q)B_W=0,7\times120=84$, \quad $\E[u_N]=-D-k_T=-210$.
\paragraph{$O_3$ — Times censuré, Post muet.}
$\E[u_T]=q(B_T-c_T)+(1-q)(\theta B_T-c_T-S_T)=27-21=6$, \quad $\E[u_W]=0$, \quad $\E[u_N]=q(-D-k_T)+(1-q)(-\theta D-k_T)=-63-49=-112$.
\paragraph{$O_5$ — procès du Post seul.}
$\E[u_T]=-[q+(1-q)\theta]R_T=-25,5$, \quad $\E[u_W]=q(B_W-c_W)+(1-q)(\theta B_W-c_W-S_W)=31,5-111,3=-79,8$, \quad $\E[u_N]=q(-D-k_W)+(1-q)(-\theta D-k_W)=-66-56=-122$.
\paragraph{Feuilles sans procès.} Aucun verdict, le gain est certain : $O_4=(B_T;\,0;\,-D)=(100;\,0;\,-200)$, $O_6=(-R_T;\,B_W;\,-D)=(-50;\,120;\,-200)$ et $O_7=(0;\,0;\,0)$.

\begin{intuition}
Dans $O_2$ et $O_3$, avec la probabilité $q$ le juge refuse de censurer le Times : celui-ci publie tout lui-même et le choix du Post ne change rien. C'est pourquoi le Post ne reçoit que $(1-q)B_W$ dans $O_2$.
\end{intuition}

\subsection{Deux lemmes de seuil}
\begin{lemme}[seuil d'un journal]\label{lem:journal}
Un journal qui publie et sera poursuivi gagne finalement avec une probabilité $\pi$. Son gain espéré, comparé à $0$ s'il se tait, vaut
\[
G(\pi)=\pi\,(B-c)+(1-\pi)\,(\theta B-c-S).
\]
$G$ est croissante en $\pi$ (pente $(1-\theta)B+S>0$), donc le journal publie si et seulement si
\[
\pi\;\ge\;\tau\;=\;\frac{c+S-\theta B}{(1-\theta)B+S}.
\]
Pour le Post, $\tau_W=\dfrac{15+180-36}{84+180}=\dfrac{159}{264}=\dfrac{53}{88}\approx60,2\,\%$. Pour le Times, $\tau_T=\dfrac{10+50-30}{70+50}=\dfrac14=25\,\%$.
\end{lemme}
\begin{proof}
$G(\pi)\ge0 \iff \pi\big[(B-c)-(\theta B-c-S)\big]\ge-(\theta B-c-S) \iff \pi\big[(1-\theta)B+S\big]\ge c+S-\theta B$.
\end{proof}

\begin{lemme}[seuil d'un procès pour Nixon]\label{lem:nixon}
Si la censure tient jusqu'au bout avec la probabilité $1-\pi$, un procès évite en moyenne le dommage $(1-\pi)(1-\theta)D$ (la part non encore publiée), pour un coût $k$. Nixon poursuit si et seulement si $(1-\pi)(1-\theta)D>k$. Ici $(1-\theta)D=140$.
\end{lemme}

\section{Induction à rebours en fonction de $q$ et $p$}
On résout chaque nœud après toutes ses suites. À chaque nœud, le joueur compare \emph{sa} composante ; on fait remonter le triplet entier (S6, diapositives 21--24). Les différences ci-dessous partent des gains espérés de la section 3.

\paragraph{$x_3$ — Nixon face au Post en relais.}
\begin{align*}
\E[u_N(O_1)]-\E[u_N(O_2)]&=(1-q)\big[(1-p)(1-\theta)D-k_W\big]>0\\
\iff p<\bar p_N&=1-\frac{k_W}{(1-\theta)D}=1-\frac{20}{140}=\frac67\approx85,7\,\%.
\end{align*}
Avec $p=0,65$ : $0,35\times140=49>20$, Nixon poursuit. La menace est \textbf{crédible} (lemme \ref{lem:nixon} avec $\pi=p$, car le Times est déjà censuré).

\paragraph{$x_2$ — le Post en relais.} Si Nixon le poursuivrait :
\begin{align*}
\E[u_W(O_1)]-\E[u_W(O_3)]&=(1-q)\big[p(B_W-c_W)+(1-p)(\theta B_W-c_W-S_W)\big]\ge0\\
\iff p\ge\tau_W&=\frac{53}{88}\approx60,2\,\%.
\end{align*}
Avec $p=0,65$ : $8,82>0$, le Post publie, avec une marge de $4,8$ points seulement. Si Nixon ne le poursuivrait pas ($p\ge\frac67$), il gagne $(1-q)B_W\ge0$ et publie toujours.

\paragraph{$x_1$ — Nixon face au Times.} Trois cas, selon ce qui suit :
\begin{itemize}[leftmargin=*,itemsep=2pt]
\item le Post publierait et serait poursuivi ($\tau_W\le p<\frac67$) :
\begin{align*}
\E[u_N(O_1)]-\E[u_N(O_4)]&=(1-q)(1-p)(1-\theta)D-k_T-(1-q)k_W>0\\
\iff q<\bar q_A(p)&=1-\frac{k_T}{(1-p)(1-\theta)D-k_W}\ ;
\end{align*}
avec $p=0,65$ : $\bar q_A=1-\frac{10}{29}=\frac{19}{29}\approx65,5\,\%$, et ici $24<34,3$, Nixon poursuit ;
\item le Post se tairait ($p<\tau_W$) : $\E[u_N(O_3)]-\E[u_N(O_4)]=(1-q)(1-\theta)D-k_T>0\iff q<1-\frac{10}{140}=\frac{13}{14}\approx92,9\,\%$ ;
\item le Post publierait sans être poursuivi ($p\ge\frac67$) : la différence vaut $-k_T<0$ ; poursuivre le Times ne sert à rien.
\end{itemize}

\paragraph{$x_5$ — Nixon face au Post seul.} Le Post seul n'a que le juge local ($\pi=q$) :
$(1-q)(1-\theta)D>k_W\iff q<\frac67\approx85,7\,\%$. Avec $q=0,3$ : $98>20$, Nixon poursuivrait.

\paragraph{$x_4$ — le Post seul.} Lemme \ref{lem:journal} avec $\pi=q$ : il publie si $q\ge\tau_W\approx60,2\,\%$. Avec $q=0,3$ : $\E[u_W(O_5)]=-79,8<0$, il se tairait. \emph{Sans le Times, rien ne sortirait.}

\paragraph{$x_0$ — le Times.} Avec Nixon qui poursuit et le Post qui relaie, le Times gagne finalement avec $\pi=q+(1-q)p=0,755$ ; lemme \ref{lem:journal} : $0,755\ge\tau_T=0,25$, donc il publie ($60,6>0$, l'alternative étant $O_7$). Sans relais, il ne compterait que sur le juge local : $\pi=q$, et il publierait encore tant que $q\ge\frac14$.

\begin{resultat}
\textbf{ENPS} : Times $=$ publier ; Post $=$ (publier en relais ; se taire seul) ; Nixon $=$ (poursuivre le Times ; poursuivre le Post ; poursuivre le Post seul). Feuille atteinte : $O_1$, gains espérés $(60,6;\ 8,82;\ -189,7)$.
\end{resultat}

\begin{cours}
L'ENPS fixe une action à chaque nœud, y compris $x_4$ et $x_5$, jamais atteints en 1971 (S6, diapositives 10--18). Une menace se juge à ce qu'elle rapporterait si on devait l'exécuter, comme la guerre dans le jeu de l'entrée (S6, diapositives 3--6) : ici, poursuivre le Post n'est crédible que si $p<\frac67$.
\end{cours}

\section{Pour quelles probabilités publie-t-on, attaque-t-on ?}

\subsection{Synthèse des seuils}
\begin{center}
\small
\begin{tabular}{@{}lll@{}}
\toprule
Décision & Condition & Valeur \\
\midrule
Le Post publie en relais & $p\ge\tau_W=\frac{c_W+S_W-\theta B_W}{(1-\theta)B_W+S_W}$ (ou $p\ge\bar p_N$) & $60,2\,\%$ \\
Le Post publie seul & $q\ge\tau_W$ (ou $q\ge\frac67$) & $60,2\,\%$ \\
Le Times publie & $q+(1-q)p\ge\tau_T$ si relais ; $q\ge\tau_T$ sinon & $25\,\%$ \\
Nixon poursuit le Post en relais & $p<1-\frac{k_W}{(1-\theta)D}$ & $85,7\,\%$ \\
Nixon poursuit le Post seul & $q<1-\frac{k_W}{(1-\theta)D}$ & $85,7\,\%$ \\
Nixon poursuit le Times (relais) & $q<1-\frac{k_T}{(1-p)(1-\theta)D-k_W}$ & $65,5\,\%$ en $p=0,65$ \\
Nixon poursuit le Times (Post muet) & $q<1-\frac{k_T}{(1-\theta)D}$ & $92,9\,\%$ \\
Deux procès possibles & $p<1-\frac{k_T+k_W}{(1-\theta)D}$ & $78,6\,\%$ \\
\bottomrule
\end{tabular}
\end{center}
Le coût total des procédures intervient ainsi : les deux procès ont lieu si $k_T+(1-q)k_W<(1-q)(1-p)(1-\theta)D$ (coût total espéré $24$ contre dommage espéré évité $34,3$) et si $k_W<(1-p)(1-\theta)D$ ($20<49$).

\subsection{La carte des équilibres}
\begin{figure}[H]
\centering
\begin{tikzpicture}
\begin{axis}[width=10.5cm,height=10.5cm,xmin=0,xmax=1,ymin=0,ymax=1,
  xlabel={$q$ : probabilité qu'un juge local refuse la censure},
  ylabel={$p$ : probabilité que la Cour suprême donne raison aux journaux},
  xtick={0,0.25,0.5,0.75,1},ytick={0,0.25,0.5,0.75,1},
  xticklabels={0,25\,\%,50\,\%,75\,\%,100\,\%},yticklabels={0,25\,\%,50\,\%,75\,\%,100\,\%},
  label style={font=\small},tick label style={font=\small},axis on top]
\fill[regnone] (axis cs:0,0) rectangle (axis cs:0.25,0.6022727272727273);
\fill[regnyt] (axis cs:0.25,0) rectangle (axis cs:0.9285714285714286,0.6022727272727273);
\fill[regnosuit] (axis cs:0.9285714285714286,0) rectangle (axis cs:1,0.6022727272727273);
\fill[regnosuit] (axis cs:0,0.6022727272727273) rectangle (axis cs:1,1);
\fill[regboth] plot coordinates {(0.0000,0.6023) (0.7197,0.6023) (0.7163,0.6053) (0.7129,0.6084) (0.7093,0.6114) (0.7056,0.6145) (0.7019,0.6176) (0.6980,0.6206) (0.6941,0.6237) (0.6900,0.6267) (0.6858,0.6298) (0.6815,0.6328) (0.6771,0.6359) (0.6726,0.6390) (0.6680,0.6420) (0.6632,0.6451) (0.6583,0.6481) (0.6532,0.6512) (0.6480,0.6542) (0.6426,0.6573) (0.6370,0.6604) (0.6313,0.6634) (0.6254,0.6665) (0.6193,0.6695) (0.6130,0.6726) (0.6064,0.6756) (0.5997,0.6787) (0.5927,0.6818) (0.5855,0.6848) (0.5780,0.6879) (0.5702,0.6909) (0.5622,0.6940) (0.5538,0.6971) (0.5451,0.7001) (0.5361,0.7032) (0.5267,0.7062) (0.5169,0.7093) (0.5067,0.7123) (0.4961,0.7154) (0.4850,0.7185) (0.4734,0.7215) (0.4612,0.7246) (0.4485,0.7276) (0.4352,0.7307) (0.4212,0.7337) (0.4065,0.7368) (0.3910,0.7399) (0.3747,0.7429) (0.3575,0.7460) (0.3393,0.7490) (0.3201,0.7521) (0.2997,0.7551) (0.2781,0.7582) (0.2551,0.7613) (0.2305,0.7643) (0.2043,0.7674) (0.1763,0.7704) (0.1462,0.7735) (0.1138,0.7765) (0.0789,0.7796) (0.0410,0.7827) (0.0000,0.7857) (0.0000,0.7857)} -- cycle;
\fill[pattern=north east lines, pattern color=green!45!black] (axis cs:0,0.8571428571428572) rectangle (axis cs:1,1);
\draw[thick] (axis cs:0,0.6022727272727273) -- (axis cs:1,0.6022727272727273);
\draw[thick] (axis cs:0.25,0) -- (axis cs:0.25,0.6022727272727273);
\draw[thick] (axis cs:0.9285714285714286,0) -- (axis cs:0.9285714285714286,0.6022727272727273);
\draw[thick] plot coordinates {(0.7197,0.6023) (0.7163,0.6053) (0.7129,0.6084) (0.7093,0.6114) (0.7056,0.6145) (0.7019,0.6176) (0.6980,0.6206) (0.6941,0.6237) (0.6900,0.6267) (0.6858,0.6298) (0.6815,0.6328) (0.6771,0.6359) (0.6726,0.6390) (0.6680,0.6420) (0.6632,0.6451) (0.6583,0.6481) (0.6532,0.6512) (0.6480,0.6542) (0.6426,0.6573) (0.6370,0.6604) (0.6313,0.6634) (0.6254,0.6665) (0.6193,0.6695) (0.6130,0.6726) (0.6064,0.6756) (0.5997,0.6787) (0.5927,0.6818) (0.5855,0.6848) (0.5780,0.6879) (0.5702,0.6909) (0.5622,0.6940) (0.5538,0.6971) (0.5451,0.7001) (0.5361,0.7032) (0.5267,0.7062) (0.5169,0.7093) (0.5067,0.7123) (0.4961,0.7154) (0.4850,0.7185) (0.4734,0.7215) (0.4612,0.7246) (0.4485,0.7276) (0.4352,0.7307) (0.4212,0.7337) (0.4065,0.7368) (0.3910,0.7399) (0.3747,0.7429) (0.3575,0.7460) (0.3393,0.7490) (0.3201,0.7521) (0.2997,0.7551) (0.2781,0.7582) (0.2551,0.7613) (0.2305,0.7643) (0.2043,0.7674) (0.1763,0.7704) (0.1462,0.7735) (0.1138,0.7765) (0.0789,0.7796) (0.0410,0.7827) (0.0000,0.7857)};
\draw[thick, dashed] (axis cs:0,0.8571428571428572) -- (axis cs:1,0.8571428571428572);
\node[circle,fill=gov,inner sep=2.2pt,label={[font=\small\bfseries]right:1971}] at (axis cs:0.3,0.65) {};
\node[font=\footnotesize,align=center] at (axis cs:0.12,0.3) {Rien\\ne sort};
\node[font=\footnotesize,align=center] at (axis cs:0.58,0.3) {Procès du Times seul\\(le Post se tait)};
\node[font=\footnotesize,align=center] at (axis cs:0.25,0.69) {Deux procès};
\node[font=\footnotesize,align=center] at (axis cs:0.72,0.75) {Sans procès};
\node[font=\footnotesize,align=center,fill=white,inner sep=1pt] at (axis cs:0.5,0.93) {Sans procès, menace non crédible};
\end{axis}
\end{tikzpicture}
\caption{Équilibre selon $(q,p)$, les autres paramètres restant ceux de référence. Frontières : $p=\tau_W\approx0,602$ ; $q=\frac14$ et $q=\frac{13}{14}$ sous cette ligne ; la courbe $q=1-\frac{10}{140(1-p)-20}$ entre $p=0,602$ et $p=\frac{11}{14}\approx0,786$ ; pointillés : $p=\frac67$.}
\end{figure}

\begin{itemize}[leftmargin=*,itemsep=1pt]
\item \textbf{Rien ne sort} : $q<\frac14$ et $p<\tau_W$. Le Post ne relaierait pas, le Times n'aurait que le juge local et ses chances passent sous son seuil ; seul, le Post n'oserait pas non plus.
\item \textbf{Procès du Times seul} : $\frac14\le q<\frac{13}{14}$ et $p<\tau_W$. Le Times publie, Nixon le poursuit, le Post se tait. La censure tient avec la probabilité $1-q$.
\item \textbf{Deux procès} (1971) : $\tau_W\le p<\frac{11}{14}$ et $q<\bar q_A(p)$.
\item \textbf{Sans procès} : partout ailleurs. Au-dessus de $p=\frac67$, la menace contre le Post n'est plus crédible.
\end{itemize}

\section{Les probabilités des issues finales}
À l'équilibre, la feuille atteinte est certaine ; seuls les verdicts restent aléatoires :
\begin{center}
\small
\begin{tabular}{@{}lll@{}}
\toprule
Régime & $\Pr(\text{tout sort})$ & $\Pr(\text{censure maintenue})$ \\
\midrule
Deux procès & $q+(1-q)p$ & $(1-q)(1-p)$ \\
Procès du Times seul & $q$ & $1-q$ \\
Sans procès & $1$ & $0$ \\
Rien ne sort & $0$ & $0$ (rien n'est publié) \\
\bottomrule
\end{tabular}
\end{center}
En 1971 : juge local favorable au Times $q=30\,\%$ ; censure puis victoire à la Cour suprême $(1-q)p=45,5\,\%$ ; censure confirmée $(1-q)(1-p)=24,5\,\%$. Total « tout sort » : $75,5\,\%$.

\begin{table}[H]
\centering
\small
\caption{Probabilité que tout sorte à l'équilibre, selon $q$ (colonnes) et $p$ (lignes). Lettre : D deux procès, T procès du Times seul, R rien ne sort, S sans procès, S* sans procès avec menace non crédible.}
\begin{tabular}{@{}lcccccc@{}}
\toprule
 & $q=0,1$ & $q=0,2$ & $q=0,3$ & $q=0,5$ & $q=0,7$ & $q=0,9$ \\
\midrule
$p=0,9$ & $100\,\%$\,{\scriptsize S*} & $100\,\%$\,{\scriptsize S*} & $100\,\%$\,{\scriptsize S*} & $100\,\%$\,{\scriptsize S*} & $100\,\%$\,{\scriptsize S*} & $100\,\%$\,{\scriptsize S*} \\
$p=0,8$ & $100\,\%$\,{\scriptsize S} & $100\,\%$\,{\scriptsize S} & $100\,\%$\,{\scriptsize S} & $100\,\%$\,{\scriptsize S} & $100\,\%$\,{\scriptsize S} & $100\,\%$\,{\scriptsize S} \\
$p=0,75$ & $77,5\,\%$\,{\scriptsize D} & $80\,\%$\,{\scriptsize D} & $82,5\,\%$\,{\scriptsize D} & $100\,\%$\,{\scriptsize S} & $100\,\%$\,{\scriptsize S} & $100\,\%$\,{\scriptsize S} \\
$p=0,7$ & $73\,\%$\,{\scriptsize D} & $76\,\%$\,{\scriptsize D} & $79\,\%$\,{\scriptsize D} & $85\,\%$\,{\scriptsize D} & $100\,\%$\,{\scriptsize S} & $100\,\%$\,{\scriptsize S} \\
$p=0,65$ & $68,5\,\%$\,{\scriptsize D} & $72\,\%$\,{\scriptsize D} & $75,5\,\%$\,{\scriptsize D} & $82,5\,\%$\,{\scriptsize D} & $100\,\%$\,{\scriptsize S} & $100\,\%$\,{\scriptsize S} \\
$p=0,6$ & $0\,\%$\,{\scriptsize R} & $0\,\%$\,{\scriptsize R} & $30\,\%$\,{\scriptsize T} & $50\,\%$\,{\scriptsize T} & $70\,\%$\,{\scriptsize T} & $90\,\%$\,{\scriptsize T} \\
$p=0,5$ & $0\,\%$\,{\scriptsize R} & $0\,\%$\,{\scriptsize R} & $30\,\%$\,{\scriptsize T} & $50\,\%$\,{\scriptsize T} & $70\,\%$\,{\scriptsize T} & $90\,\%$\,{\scriptsize T} \\
$p=0,3$ & $0\,\%$\,{\scriptsize R} & $0\,\%$\,{\scriptsize R} & $30\,\%$\,{\scriptsize T} & $50\,\%$\,{\scriptsize T} & $70\,\%$\,{\scriptsize T} & $90\,\%$\,{\scriptsize T} \\
\bottomrule
\end{tabular}
\end{table}

\begin{intuition}
La ligne $p=0,6$ tombe juste sous le seuil du Post ($60,2\,\%$) : il se tait et la probabilité que tout sorte chute de $75\,\%$ environ à $q$. Une petite variation de la croyance sur la Cour suprême fait basculer l'issue : c'est ce qui rendait la décision du Post si difficile.
\end{intuition}

\section{Statique comparative : l'effet de $q$ et de $p$}
\begin{itemize}[leftmargin=*,itemsep=3pt]
\item \textbf{Le gain du Post en relais} croît avec $p$ : $\dfrac{\partial\,\E[u_W(O_1)]}{\partial p}=(1-q)\big[(1-\theta)B_W+S_W\big]=0,7\times264=184,8$. Un point de pourcentage de $p$ en plus vaut $1,85$ pour le Post.
\item \textbf{L'intérêt de Nixon à poursuivre le Post} décroît avec $p$ : $\dfrac{\partial}{\partial p}\big[(1-p)(1-\theta)D-k_W\big]=-(1-\theta)D=-140$.
\item \textbf{L'intérêt de Nixon à poursuivre le Times} décroît avec $q$ : $\dfrac{\partial}{\partial q}\big[(1-q)\big((1-p)(1-\theta)D-k_W\big)-k_T\big]=-\big[(1-p)(1-\theta)D-k_W\big]=-29<0$.
\item \textbf{La probabilité que tout sorte} en régime « deux procès », $\pi=q+(1-q)p$, croît avec $q$ ($\partial\pi/\partial q=1-p=0,35$) et avec $p$ ($\partial\pi/\partial p=1-q=0,7$).
\item \textbf{Marges de 1971} : $p=0,65$ dépasse le seuil du Post de $4,8$ points et reste sous le seuil de Nixon de $20,7$ points ; $q=0,3$ reste $35,5$ points sous le seuil de poursuite du Times.
\end{itemize}

\section{Cas chiffrés}
Chaque ligne change seulement les valeurs indiquées ; les autres restent celles de référence. Les gains sont ceux du nœud initial (Times ; Post ; Nixon).
\begin{center}
\scriptsize
\begin{tabularx}{\textwidth}{@{}>{\raggedright\arraybackslash}p{3.1cm}>{\raggedright\arraybackslash}X>{\raggedright\arraybackslash}p{2.6cm}cll@{}}
\toprule
Scénario & Changement & Régime & Feuille & Gains espérés & $\Pr(\text{tout sort})$ \\
\midrule
1971 · scénario de référence & référence & Deux procès & $O_{1}$ & $(60,6;\,8,82;\,-189,7)$ & $75,5\,\%$ \\
Le Post prend peur & $p=0,5$ & Procès du Times seul & $O_{3}$ & $(6;\,0;\,-112)$ & $30\,\%$ \\
Sanction du Post à 250 & $S_W=250$ & Procès du Times seul & $O_{3}$ & $(6;\,0;\,-112)$ & $30\,\%$ \\
Le Times renonce & $c_T=80$ & Rien ne sort & $O_{7}$ & $(0;\,0;\,0)$ & $0\,\%$ \\
Tribunaux hostiles & $q=0,2$, $p=0,5$ & Rien ne sort & $O_{7}$ & $(0;\,0;\,0)$ & $0\,\%$ \\
Procès du Post trop cher & $k_W=40$ & Sans procès & $O_{4}$ & $(100;\,0;\,-200)$ & $100\,\%$ \\
Menace non crédible & $k_W=60$ & Sans procès (menace non crédible) & $O_{4}$ & $(100;\,0;\,-200)$ & $100\,\%$ \\
Cour suprême acquise & $p=0,9$ & Sans procès (menace non crédible) & $O_{4}$ & $(100;\,0;\,-200)$ & $100\,\%$ \\
Juge local favorable & $q=0,8$ & Sans procès & $O_{4}$ & $(100;\,0;\,-200)$ & $100\,\%$ \\
Publier pour ne pas être doublé & $q=0,65$, $p=0,5$, $c_T=30$, $S_T=150$ & Procès du Times seul & $O_{3}$ & $(-7;\,0;\,-161)$ & $65\,\%$ \\
Menace crue à tort & $p=0,5$, $k_W=80$ & Sans procès (menace non crédible) & $O_{4}$ & $(100;\,0;\,-200)$ & $100\,\%$ \\
\bottomrule
\end{tabularx}
\end{center}

\paragraph{Équilibre de Nash non parfait ($p=0,5$, $k_W=80$).} Poursuivre le Post rapporterait à Nixon $-10-80-200(0,5+0,5\times0,3)=-220<-210$ : la menace n'est pas crédible. Pourtant, le profil où Nixon « menace » de poursuivre le Post et où le Post se tait est un équilibre de Nash : la menace ne coûte rien tant qu'elle n'est pas mise à l'épreuve, et le secret tient avec la probabilité $1-q=70\,\%$. L'ENPS l'élimine : le Post publierait, Nixon ne poursuivrait personne, et tout sort ($O_4$). Le recensement des 64 profils donne 4 équilibres de Nash dans le jeu de référence (un seul ENPS) et 12 dans ce scénario.

\section{Conventions et limites}
\begin{itemize}[leftmargin=*,itemsep=1pt]
\item À égalité de gains, un journal publie et Nixon ne poursuit pas.
\item Les verdicts sont supposés indépendants et leurs probabilités connues de tous (croyances communes).
\item La Cour suprême n'est saisie en urgence que si deux journaux sont poursuivis ; sinon la censure locale tient.
\item Les valeurs sont un calibrage illustratif en points d'utilité, pas des estimations historiques. L'usage d'espérances suppose que ces points mesurent l'intensité des préférences.
\end{itemize}

\end{document}
